\section{问题重述}

\subsection{问题背景}

随着"双碳"战略深入推进，新能源装机规模持续快速增长，但新能源大规模并网消纳难的问题日益凸显，成为制约能源绿色转型的关键瓶颈\cite{ndrc2024,liu2015}。绿电直连项目通过将新能源发电端与用电负荷直接连接，为促进新能源就近消纳提供了新路径。其中，"绿电—绿氢—绿氨"一体化电氢氨化工园区，利用可再生能源电解水制氢再将氢转化为氨，既解决了绿电就地消纳问题，又为化工行业提供了深度脱碳的可行路径\cite{wang2024}。因此，研究绿电直连型电氢氨园区的优化运行问题，对推动新能源消纳和化工行业绿色转型具有重要的理论价值和现实意义。

本文以某典型绿电直连制氨项目为研究对象，园区包含风电(40MW)、光伏(64MW)、碱性电解槽、质子交换膜电解槽、合成氨装置、常规负荷及联网线路。给定24种风、光出力组合场景及分时电价等参数，研究园区的优化运行问题。

\subsection{问题重述}

\textbf{问题一：典型风光场景下的运行指标分析。}给定典型日风光出力及常规负荷曲线，电解槽与合成氨装置满负荷连续运行(36吨/日)。要求根据功率平衡计算购售电功率并绘制曲线，计算绿电直连三项指标及吨氨成本，分析达标情况及原因。

\textbf{问题二：基于离散制氨调节的运行优化。}制氨产能增至72吨/日，产量从72吨/日按9吨递减至36吨/日，电氢氨装置仅全额开机或停机。要求在典型场景下给出各产量成本最低的生产时段安排及绿电指标，找出最优日产量；在24种风光场景下统计分析各指标的分布特征，按全满足、部分满足、全不满足三类统计全年绿电指标状况，绘制成本分布曲线并计算全年总吨氨成本。

\textbf{问题三：基于连续制氨调节的运行分析。}产能72吨/日，产量从72吨/日递减至36吨/日，电氢氨装置功率连续可调(下限10\%)。要求在24种风光场景下计算最优功率调度方案及相应指标和成本，按三类情况统计全年绿电指标，分析各场景运行原因，并与问题二结果对比分析变化情况。

\textbf{问题四：离网运行分析及储能配置研究。}园区离网运行，制氨产能72吨/日，功率连续可调(下限10\%)。要求在尽限利用风光条件下计算各场景产量、成本及产能利用率，分析能源自给状况；给出最大弃电场景的储能最优配置方案；基于相同产量需求对比离网和联网两种运行模式的经济性。

\textbf{问题五：政策影响分析与建议。}分析绿电园区容量渗透率显著提高后对电力系统运行的影响（至少三种利弊），提出推动绿电直连园区发展的政策建议并给出论证依据。
